% Generated by Claude Fable 5.1 (Anthropic), 2026-09-14 — Fig. 2 caption, panel order (a) graph, (b) decision tree, (c) layers, (d) GPU
\caption{The FISH approach. (a)~The fully expanded FISH graph: the containment
hierarchy from the host down to the individual callbacks, with colour marking
the namespace a vertex belongs to and the edges being the containment relation.
(b)~How the kind of an edge is decided: a shared topic or service, known from
initialisation, gives a $\mathit{sub}$ or $\mathit{dep}$ edge; the remaining two
are read off execution alone. (c)~How edges move up the layers. The edges marked
\textcircled{\scriptsize 2} carry the very same message, drawn again at every
layer above. The two marked \textcircled{\scriptsize 3} are one output of a join:
either member callback may publish it, whichever completes the set, so the
message leaves two callbacks and two entities; at \textcircled{\scriptsize 4},
where both endpoints are nodes, they merge into one edge. Edges whose two
endpoints sit inside one node, $\mathit{polled\text{-}msg}$ and
$\mathit{join}$, are contained by it: at the node layer they lie inside the
vertex, and expanding the node brings them back. (d)~For a GPU-submitting
callback FISH keeps every execution as a typed DAG, groups the executions into
shapes, and folds the shapes into one conditional graph: a skeleton common to
all executions plus the branches, each with the share of executions that take
it.}
